\chapter{Herleitung der Methode der kleinsten Quadrate}\label{sec:appendixKleinsteQuadrate}

\section*{Definition der Zielfunktion}
Das Ziel der linearen Regression besteht darin, die bestmögliche Gerade $y = mx + q$ zu finden, die die Summe der quadrierten Residuen zwischen den vorhergesagten Werten und den tatsächlichen Datenpunkten minimiert. Die Summe der quadrierten Abstände zwischen den Datenpunkten und der Geraden wird definiert als:
$$
S = \sum_{i=1}^n (y_i - (mx_i + q))^2
$$

\section*{Partielle Ableitungen}
Um $S$ zu minimieren, nehmen wir die partiellen Ableitungen bezüglich $m$ und $q$ und setzen diese gleich Null.

\subsection*{Partielle Ableitung nach $m$}
$$
\frac{\partial S}{\partial m} = \frac{\partial}{\partial m} \sum_{i=1}^n (y_i - mx_i - q)^2
$$
$$
= \sum_{i=1}^n 2(y_i - mx_i - q)(-x_i)
$$
$$
= -2 \sum_{i=1}^n x_i(y_i - mx_i - q)
$$
Auf Null setzen für die Minimierung:
$$
\sum x_iy_i - m\sum x_i^2 - q\sum x_i = 0
$$
$$
\sum x_iy_i = m\sum x_i^2 + q\sum x_i \quad \text{(Gleichung 1)}
$$

\subsection*{Partielle Ableitung nach $q$}
$$
\frac{\partial S}{\partial q} = \frac{\partial}{\partial q} \sum_{i=1}^n (y_i - mx_i - q)^2
$$
$$
= \sum_{i=1}^n 2(y_i - mx_i - q)(-1)
$$
$$
= -2 \sum_{i=1}^n (y_i - mx_i - q)
$$
Auf Null setzen für die Minimierung:
$$
\sum y_i - m\sum x_i - nq = 0
$$
$$
\sum y_i = m\sum x_i + nq \quad \text{(Gleichung 2)}
$$

\section*{Lösung der Gleichungen}
Aus Gleichung 2:
$$
q = \frac{\sum y_i - m\sum x_i}{n} \quad \text{(Gleichung 3)}
$$
Einsetzen der Gleichung 3 in Gleichung 1:
$$
\sum x_iy_i = m\sum x_i^2 + \left(\frac{\sum y_i - m\sum x_i}{n}\right)\sum x_i
$$
$$
\sum x_iy_i = m\sum x_i^2 + \frac{(\sum y_i)\sum x_i}{n} - \frac{m(\sum x_i)^2}{n}
$$
Auflösen nach $m$:
$$
m = \frac{n\sum x_iy_i - \sum x_i \sum y_i}{n\sum x_i^2 - (\sum x_i)^2}
$$
Bestimmen von $q$ über Einsetzen von $m$ in Gleichung 3:
$$
q = \frac{\sum y_i - m\sum x_i}{n}
$$